% Printed 262; principal PDF 269; Google PDF 273.
\par\clearpage
\null\thispagestyle{empty}

% Printed 263; principal PDF 270; Google PDF 274.
\par\clearpage
\thispagestyle{empty}
\vspace*{2.8\baselineskip}
{\centering\Large\bfseries NACHLASS UND BRIEFWECHSEL.\par}
\vspace{3.3\baselineskip}
{\centering
[1.]\par\smallskip
[Stationsausgleichung für] Zeven aus sämtlichen Messungen\\
von 1824 und 1825 (ohne die vom 4. und 5. August).\par\smallskip
[Annahme für die Azimute.]\par}
\[
\begin{array}{@{}r|l|r@{\;}r@{\;}r@{}l@{}}
 1&\text{Steinberg}&1^\circ&17'&45''{,}262&+a\\
 &\text{Brüttendorf, Zentr.}&17&21&48{,}275&\\
 2&\text{\hbox to 8.5em{\quad—\hfil Strich}}&17&21&58{,}264&+f\\
 &\text{Bremen, Zentr.}&53&26&10{,}871&\\
 3&\text{\hbox to 8.5em{\quad—\hfil Heliotrop}}&53&26&13{,}375&+b\\
 4&\text{Brillit}&124&13&47{,}128&+c\\
 5&\text{Selsingen}&152&47&56{,}646&+d\\
 6&\text{Litberg}&246&0&49{,}110&+e\\
 7&\text{Wilsede}&288&22&42{,}653&+f
\end{array}
\]
{\small
\[
\begin{array}{@{}c@{\quad}r@{\;}c@{\;}r@{\;}c@{\;}r@{\quad}r@{\;}r@{\;}r@{}}
 \Bigl[\,\substack{\text{Beobachtete}\\\text{Winkel}}&
 \multicolumn{5}{c}{\text{Anzahl der Repetitionen}}&
 \multicolumn{3}{c@{}}{\text{Beobachtungswerte}\,\Bigr]}\\[-2pt]
 &&&\multicolumn{3}{c}{\text{Kreis}}&&&\\[-3pt]
 &&&\mathit{l.}&&\mathit{r.}&&&\\
 1.3&57&=&35&+&22&52^\circ&8'&27''{,}798\\
 1.4&68&&40&&28&122&56&1{,}603\\
 1.5&9&&2&&7&151&30&11{,}417\\
 3.4&132&&70&&62&70&47&33{,}506\\
 3.5&16&&2&&14&99&21&44{,}047.\\
 4.5&12&&\ast&&12&28&34&8{,}541\\
 4.7&48&&30&&18&164&8&54{,}719\\
 6.7&45&&15&&30&42&21&53{,}578\\
 6.1&53&&29&&24&115&16&56{,}075
\end{array}
\]
}

% Printed 264; principal PDF 271; Google PDF 275.
\par\clearpage
\noindent\makebox[0pt][l]{\fontsize{13}{16}\selectfont 264}\hfill{\small\textbf{NACHLASS.}}\hfill\null\par\medskip
{\small
\[
\begin{array}{@{}c@{\quad}r@{\;}c@{\;}r@{\;}c@{\;}r@{\quad}r@{\;}r@{\;}r@{}}
 \Bigl[\,\substack{\text{Beobachtete}\\\text{Winkel}}&
 \multicolumn{5}{c}{\text{Anzahl der Repetitionen}}&
 \multicolumn{3}{c@{}}{\text{Beobachtungswerte}\,\Bigr]}\\[-2pt]
 &&&\multicolumn{3}{c}{\text{Kreis}}&&&\\[-3pt]
 &&&\mathit{l.}&&\mathit{r.}&&&\\
 6.3&2&=&2&+&\ast&167^\circ&25'&25''{,}500\\
 7.1&27&&5&&22&72&55&1{,}408\\
 7.2&20&&10&&10&88&59&15{,}611\\
 7.3&3\mathord{\cdot}&&3&&\ast&125&3&29{,}167\\
 3.4+3.1&\frac12&&\mathord{\raisebox{.4ex}{.}}\ast&&\frac12&18&39&5{,}500.
\end{array}
\]
}

{\centering [Die Ausdrücke für die Fehler der Beobachtungswerte lauten dann:]\par}
\[
\begin{array}{@{}r@{\;}l@{}}
 +0''{,}315&-a+b\\
 +0{,}263&-a+c\\
 -0{,}033&-a+d\\
 +0{,}247&-b+c\\
 -0{,}776&-b+d\\
 +0{,}977&-c+d\\
 +0{,}806&-c+f\\
 -0{,}035&-e+f\\
 +0{,}077&+a-e\\
 -1{,}235&+b-e\\
 +1{,}201&+a-f\\
 \makebox[2.4em][c]{\(\ast\)}&\quad\ast\\
 +1{,}555&+b-f\\
 +0{,}140&+a-2b+c.
\end{array}
\]

[Wenn diesen Ausdrücken die Repetitionszahlen für die zugehörigen Winkel als Gewichte zugeordnet werden, so ergeben sich die folgenden Normalgleichungen:]
\[
\begin{array}{@{}r@{\;}r@{\;}r@{\;}r@{\;}r@{\;}r@{\;}c@{\;}c@{}}
 [0=\phantom{]}&+1{,}036&+214{,}5a&-58b&-67{,}5c&-9d&-53e&-27f\\
 0=\phantom{]}&-0{,}178&-58a&+212b&-133c&-16d&-2e&-3f\\
 0=\phantom{]}&+0{,}146&-67{,}5a&-133b&+260{,}5c&-12d&\ast&-48f\\
 0=\phantom{]}&-0{,}989&-9a&-16b&-12c&+37d&\ast&\ast\\
 0=\phantom{]}&-0{,}036&-53a&-2b&\ast&\ast&+100e&-45f\\
 0=]&+0{,}021&-27a&-3b&-48c&\ast&-45e&+123f.
\end{array}
\]

% Printed 265; principal PDF 272; Google PDF 276.
\par\clearpage
\noindent\null\hfill{\small\textbf{STATIONSAUSGLEICHUNGEN.}}\hfill\makebox[0pt][r]{\fontsize{13}{16}\selectfont 265}\par\medskip

[Zur Lösung dient die folgende Tabelle, in der die Werte der Konstanten und der Unbekannten in Einheiten der 3. Dezimalstelle der Bogensekunde angegeben sind.]
\[
\begin{array}{c|c|c|c|c|c}
 d=+27&a=-4&b=+2&e=-2&f=-2&e=-1\\
 \hline
 +793&-65&-181&-75&-21&+32\\
 -610&-378&+46&+50&+56&+58\\
 -178&+92&-174&-174&-78&-78\\
 +10&+46&+14&+14&+14&+14\\
 -36&+176&+172&-28&+62&-38\\
 +21&+129&+123&+213&-33&+12.
\end{array}
\]
\[
\begin{array}{r|r}
 a&-4\\
 b&+2\\
 c&0\\
 d&+27\\
 e&-3\\
 f&-2.
\end{array}
\]

\medskip
{\centering
[2.]\par\smallskip
[Stationsausgleichung für] Brillit.\par\smallskip
[Annahme für die Azimute.]\par}
\[
\begin{array}{@{}r|r@{\;}r@{\;}r@{}l@{\quad}l@{}}
 1&19^\circ&41'&57''{,}982&+a&\text{Bremen, Heliotrop}\\
 2&49&19&52{,}539&+b&\text{Garlste}\\
 3&124&0&37{,}967&+c&\text{Bremerlehe, Heliotrop}\\
 4&304&13&51{,}679&+d&\text{Zeven, Heliotrop.}
\end{array}
\]
{\centering[Beobachtungstabelle.]\par}
{\small
\[
\begin{array}{c|c|rr|rr}
 \Bigl[\,\substack{\text{Gemessene}\\\text{Winkel}}&
 \substack{\text{Anzahl der}\\\text{Repetitionen}}\,\Bigr]&
 \multicolumn{2}{c|}{\substack{\text{[Kreis]}\\\mathit{l.}\qquad\mathit{r.}}}&
 \text{dir.}&\text{suppl.}\\
 \hline
 1.2&42&22&20&21&21\\
 2.3&48&24&24&24&24\\
 4.1&44&23&21&23&21\\
 4.2&40&20&20&20&20\\
 4.3&24&7&17&13&11.
\end{array}
\]
}
\noindent{\small IX.}\hfill{\small 34}\hspace*{.12\linewidth}\par

% Printed 266; principal PDF 273; Google PDF 277.
\par\clearpage
\noindent\makebox[0pt][l]{\fontsize{13}{16}\selectfont 266}\hfill{\small\textbf{NACHLASS.}}\hfill\null\par\medskip

Berücksichtigt man die konstante Korrektur der Winkel \(=x'\), ergibt sich folgende Tabelle:
\[
\begin{array}{@{}c@{\quad}c@{\quad}r@{\;}r@{\;}r@{}r@{}l@{}}
 \bigl[\,\text{Winkel}&\text{Repetitionszahl}&
 \multicolumn{4}{c@{}}{\text{Beobachtungswert}\,\bigr]}&\\
 1.2&21&29^\circ&37'&53''{,}381&+x'&\\
 &21&&&55{,}845&-x'&\\
 2.3&24&74&40&44{,}844&+x'&\\
 &24&&&46{,}021&-x'&\\
 4.1&21&75&28&5{,}631&+x'&\\
 &20&&&7{,}050&-x'&\\
 &3&&&6{,}333&+\frac13x'&\text{[*)]}\\
 4.2&15&105&6&0{,}317&+x'&\\
 &15&&&1{,}533&-x'&\\
 &10&&&0{,}450&\ast&\text{[**)]}\\
 4.3&13&179&46&45{,}962&+x'&\\
 &11&&&46{,}522&-x'.&
\end{array}
\]

[Wird] \(x'=+0''{,}723+x\) [gesetzt, so lauten die Ausdrücke für die Beobachtungsfehler der Winkel:]
\[
\begin{array}{@{}r@{\;}l@{\;}r@{}}
 +0''{,}453&-a+b&-x\\
 -0{,}565&-a+b&+x\\
 -0{,}139&-b+c&-x\\
 +0{,}130&-b+c&+x\\
 -0{,}051&+a-d&-x\\
 -0{,}024&+a-d&+x\\
 -0{,}271&+a-d&-\frac13x\\
 -0{,}180&+b-d&-x\\
 +0{,}050&+b-d&+x\\
 +0{,}410&+b-d&\ast\\
 -0{,}397&+c-d&-x\\
 +0{,}489&+c-d&+x,
\end{array}
\]
\noindent [aus denen, wenn die zugehörigen Repetitionszahlen als Gewichte angesetzt

\par\noindent\rule{4cm}{.4pt}\par
{\footnotesize
\noindent\hspace*{\parindent}[*) Im Beobachtungsbuch steht zu dieser Messung zusätzlich:] »1. und 3. [Repetition] direkt, 2. Supplement.«\par
\noindent\hspace*{\parindent}[**) Zu dieser Messung steht im Beobachtungsbuch:] »abwechselnd dir. und Suppl.; 1. 3. 5. 7. 9. direkt, 2. 4. 6. 8. 10. Supplement.«\par}

% Printed 267; principal PDF 274; Google PDF 278.
\par\clearpage
\noindent\null\hfill{\small\textbf{STATIONSAUSGLEICHUNGEN.}}\hfill\makebox[0pt][r]{\fontsize{13}{16}\selectfont 267}\par\medskip

\noindent werden, sich die folgenden Normalgleichungen ergeben:]
\[
\begin{array}{@{}r@{\;}r@{\;}r@{\;}c@{\;}c@{\;}c@{\;}c@{}}
 [0=\phantom{]}&-0{,}070&+185\frac13x&-2a&\ast&-2c&+4d\\
 0=\phantom{]}&-0{,}012&-2x&+86a&-42b&\ast&-44d\\
 0=\phantom{]}&+0{,}014&\ast&-42a&+130b&-48c&-40d\\
 0=\phantom{]}&+0{,}002&-2x&\ast&-48b&+72c&-24d\\
 0=]&-0{,}004&+4x&-44a&-40b&-24c&+108d.
\end{array}
\]

[Ihre Lösung beeinflusst die dritte Dezimalstelle der Bogensekunde in den angenommenen Werten nicht mehr; damit sind diese zugleich die Ausgleichungswerte.]

Summe der [mit ihren Repetitionszahlen multiplizierten] Quadrate [der Fehler]: \(19{,}052630.\)
\[
\begin{array}{@{}r@{\quad}r@{\;=\;}l@{}}
 \text{[Mittleres Fehlerquadrat:]}&
 \frac18\,.\,19{,}052630&2{,}381579.\\
 \text{[Mittlerer Fehler der Gewichtseinheit:]}&
 \surd 2{,}381579&\pm1''{,}543236.
\end{array}
\]

[Nach der Ausgleichung ist das] Gewicht von \(x'=+0''{,}723+x\) gleich \(185{,}1624\); [daher]
\[
\begin{array}{@{}l@{\;}l@{\;=\;}l@{}}
 \text{E.\ m.}&\left[=\frac{1{,}543}{\surd 185{,}16}\right]&\pm0''{,}113\\
 \text{E.\ pr.}&[=\pm0{,}113\,.\,0{,}674\ldots]&\pm0{,}076.
\end{array}
\]

Bei den neuen Messungen in Zeven ergibt sich
\[
\begin{gathered}
 \GaussRightBrace{\begin{array}{@{}l@{\;}r@{\;}c@{\;}l@{\;}r@{}}
 \text{Bremen—Brillit}&25&\text{[Rep.]}&\text{direkt}&70^\circ\,47'\,32''{,}810\\
 &25&\text{»}&\text{Supplement}&34{,}510
 \end{array}}\;50\text{ [Rep.]}\,.\,.\;70^\circ\,47'\,33''{,}660\\[3pt]
 x=0''{,}850.
\end{gathered}
\]

\medskip
{\centering
[3.]\par\smallskip
Wilsede\par\smallskip
aus sämtlichen Messungen von 1822 und 1824.\par\smallskip
[Annahme.]\par}
\[
\begin{array}{@{}r|l|r@{\;}r@{\;}r@{}l@{}}
 1&\text{Falkenberg}&7^\circ&51'&9''{,}430&+a\\
 2&\text{Elmhorst}&46&31&36{,}382&+b\\
 3&\text{Steinberg}&67&11&27{,}430&+c\\
 4&\text{Bottel}&72&0&39{,}776&+d\\
 5&\text{Bullerberg}&89&5&3{,}816&+e
\end{array}
\]
\noindent\hfill{\small 34\textsuperscript{*}}\hspace*{.12\linewidth}\par

% Printed 268; principal PDF 275; Google PDF 279.
\par\clearpage
\noindent\makebox[0pt][l]{\fontsize{13}{16}\selectfont 268}\hfill{\small\textbf{NACHLASS.}}\hfill\null\par\medskip

\[
\begin{array}{@{}r|l|r@{\;}r@{\;}r@{}l@{}}
 6&\text{Brüttendorf}&103^\circ&40'&9''{,}194&+f\\
 7&\text{Zeven}&108&22&39{,}286&+g\\
 8&\text{Litberg}&138&28&9{,}283&+h\\
 9&\text{Hamburg}&183&29&1{,}167&+i\\
 10&\text{Syk}&204&28&34{,}438&+k\\
 11&\text{Hohenhorn}&219&29&33{,}893&+l\\
 12&\text{Lüneburg}&253&20&5{,}387&+m\\
 13&\text{Nindorf}&275&29&13{,}209&+n\\
 14&\text{Timpenberg}&283&5&9{,}299&+o\\
 15&\text{Wulfsode}&298&29&57{,}926&+p\\
 16&\text{Breithorn}&330&3&15{,}798&+q\\
 17&\text{Hauselberg}&334&25&35{,}152&+r.
\end{array}
\]

\[
\begin{array}{@{}c|c|r@{\;}r@{\;}r|l@{}}
 \Bigl[\,\substack{\text{Beobachtete}\\\text{Winkel}}&
 \text{Gewichte}&\multicolumn{3}{c|}{\text{Beobachtungswerte}}&
 \multicolumn{1}{c@{}}{\text{Fehlerausdrücke}\,\Bigr]}\\
 \hline
 1.5&7&81^\circ&13'&53''{,}429&+0''{,}957-a+e\\
 1.6&20&95&49&0{,}287&-0{,}523-a+f\\
 2.3&40&20&39&50{,}369&+0{,}679-b+c\\
 2.6&8&57&8&33{,}469&-0{,}657-b+f\\
 2.7&29&61&51&3{,}103&\mathord{\raisebox{.55ex}{.}}-0{,}199-b+g\\
 2.8&68&91&56&33{,}140&-0{,}239-b+h\\
 2.9&3&136&57&24{,}167&+0{,}618-b+i\\
 3.6&10&36&28&41{,}225&+0{,}539-c+f\\
 3.7&11&41&11&10{,}773&+1{,}083-c+g\\
 3.8&11&71&16&42{,}455&-0{,}602-c+h\\
 3.9&12&116&17&32{,}417&+1{,}320-c+i\\
 4.5&15&17&4&23{,}800&+0{,}240-d+e\\
 4.6&20&31&39&29{,}687&-0{,}269-d+f\\
 4.7&5&36&21&59{,}300&+0{,}210-d+g\\
 5.6&20&14&35&4{,}863&+0{,}515-e+f\\
 6.7&18&4&42&30{,}458&-0{,}366-f+g\\
 6.8&50&34&48&0{,}125&-0{,}036-f+h\\
 7.8&41&30&5&30{,}067&-0{,}070-g+h\\
 8.9&30&45&0&52{,}592&-0{,}708-h+i
\end{array}
\]

% Printed 269; principal PDF 276; Google PDF 280.
\par\clearpage
\noindent\null\hfill{\small\textbf{STATIONSAUSGLEICHUNGEN.}}\hfill\makebox[0pt][r]{\fontsize{13}{16}\selectfont 269}\par\medskip

\[
\begin{array}{@{}c|c|r@{\;}r@{\;}r|l@{}}
 \Bigl[\,\substack{\text{Beobachtete}\\\text{Winkel}}&
 \text{Gewichte}&\multicolumn{3}{c|}{\text{Beobachtungswerte}}&
 \multicolumn{1}{c@{}}{\text{Fehlerausdrücke}\,\Bigr]}\\
 \hline
 8.12&10&114^\circ&51'&56''{,}275&-0''{,}171-h+m\\
 9.10&30&\mathord{\raisebox{.55ex}{.}}20&59&33{,}575&-0{,}304-i+k\\
 9.11&50&36&0&32{,}745&-0{,}019-i+l\\
 9.12&58&69&51&3{,}840&+0{,}380-i+m\\
 9.17&20&150&56&34{,}762&-0{,}777-i+r\\
 10.11&20&15&0&59{,}912&-0{,}457-k+l\\
 11.17&21&114&56&1{,}738&-0{,}479-l+r\\
 12.13&30&22&9&7{,}983&-0{,}161-m+n\\
 12.14&15&29&45&3{,}917&-0{,}005-m+o\\
 12.15&2&45&9&52{,}875&-0{,}336-m+p\\
 12.17&27&81&5&29{,}213&+0{,}552-m+r\\
 12.1&20&114&31&3{,}413&+0{,}630-m+a\\
 12.4&12&178&40&34{,}521&-0{,}132-m+d\\
 13.14&4&7&35&56{,}812&-0{,}722-n+o\\
 13.16&20&54&34&2{,}687&-0{,}098-n+q\\
 14.15&72&15&24&48{,}837&-0{,}210-o+p\\
 14.16&52&46&58&6{,}423&+0{,}076-o+q\\
 14.17&21&51&20&25{,}464&+0{,}389-o+r\\
 15.16&54&31&33&17{,}733&+0{,}139-p+q\\
 15.17&52&35&55&37{,}356&-0{,}130-p+r\\
 15.1&38&69&21&11{,}941&-0{,}437-p+a\\
 16.17&10&4&22&18{,}400&+0{,}954-q+r\\
 17.1&55&33&25&34{,}273&+0{,}005-r+a\\
 (2.3)+(2.7)&\frac12&82&30&55{,}250&-1{,}298-2b+c+g\\[2pt]
 (2.7)+(6.7)&\frac12&66&33&34{,}000&-1{,}004-b-f+2g\\[2pt]
 (6.8)+9(7.8)&\frac{1}{10}&305&37&28{,}000&+2{,}062-f-9g+10h\\[2pt]
 (6.4)+3(6.8)&\frac14&72&44&27{,}500&+3{,}349+d-4f+3h.
\end{array}
\]

% Printed 270; principal PDF 277; Google PDF 281.
\par\clearpage
\noindent\makebox[0pt][l]{\fontsize{13}{16}\selectfont 270}\hfill{\small\textbf{NACHLASS.}}\hfill\null\par\medskip

\noindent\hfil
\rotatebox{-90}{%
\begin{minipage}{.88\textheight}
[Unter Berücksichtigung der zugehörigen Gewichte entsprechen den Fehlerausdrücken die folgenden Normalgleichungen:]
\begin{center}
\resizebox{\linewidth}{!}{%
\(\begin{array}{@{}rr*{17}{@{\;}r}@{}}
 [0=\phantom{]}&+0{,}030&+140a&\ast&\ast&\ast&-7e&-20f&\ast&\ast&\ast&\ast&\ast&-20m&\ast&\ast&-38p&\ast&-55r\\
 0=\phantom{]}&+0{,}065&\ast&+150{,}5b&-41c&\ast&\ast&-7{,}5f&-31g&-68h&-3i&\ast&\ast&\ast&\ast&\ast&\ast&\ast&\ast\\
 0=\phantom{]}&-0{,}010&\ast&-41b&+84{,}5c&\ast&\ast&-10f&-10{,}5g&-11h&-12i&\ast&\ast&\ast&\ast&\ast&\ast&\ast&\ast\\
 0=\phantom{]}&-0{,}016.75&\ast&\ast&\ast&+52{,}25d&-15e&-21f&-5g&+0{,}75h&\ast&\ast&\ast&-12m&\ast&\ast&\ast&\ast&\ast\\
 0=\phantom{]}&-0{,}001&-7a&\ast&\ast&-15d&+42e&-20f&\ast&\ast&\ast&\ast&\ast&\ast&\ast&\ast&\ast&\ast&\ast\\
 0=\phantom{]}&-0{,}071.2&-20a&-7{,}5b&-10c&-21d&-20e&+150{,}6f&-18{,}1g&-54h&\ast&\ast&\ast&\ast&\ast&\ast&\ast&\ast&\ast\\
 0=\phantom{]}&-0{,}034.8&\ast&-31b&-10{,}5c&-5d&\ast&-18{,}1f&+114{,}6g&-50h&\ast&\ast&\ast&\ast&\ast&\ast&\ast&\ast&\ast\\
 0=\phantom{]}&-0{,}020.25&\ast&-68b&-11c&+0{,}75d&\ast&-54f&-50g&+222{,}25h&-30i&\ast&\ast&-10m&\ast&\ast&\ast&\ast&\ast\\
 0=\phantom{]}&+0{,}024&\ast&-3b&-12c&\ast&\ast&\ast&\ast&-30h&+203i&-30k&-50l&-58m&\ast&\ast&\ast&\ast&-20r\\
 0=\phantom{]}&+0{,}020&\ast&\ast&\ast&\ast&\ast&\ast&\ast&\ast&-30i&+50k&-20l&\ast&\ast&\ast&\ast&\ast&\ast\\
 0=\phantom{]}&-0{,}031&\ast&\ast&\ast&\ast&\ast&\ast&\ast&\ast&-50i&-20k&+91l&\ast&\ast&\ast&\ast&\ast&-21r\\
 0=\phantom{]}&-0{,}013&-20a&\ast&\ast&-12d&\ast&\ast&\ast&-10h&-58i&\ast&\ast&+174m&-30n&-15o&-2p&\ast&-27r\\
 0=\phantom{]}&+0{,}018&\ast&\ast&\ast&\ast&\ast&\ast&\ast&\ast&\ast&\ast&\ast&-30m&+54n&-4o&\ast&-20q&\ast\\
 0=\phantom{]}&+0{,}036&\ast&\ast&\ast&\ast&\ast&\ast&\ast&\ast&\ast&\ast&\ast&-15m&-4n&+164o&-72p&-52q&-21r\\
 0=\phantom{]}&+0{,}068&-38a&\ast&\ast&\ast&\ast&\ast&\ast&\ast&\ast&\ast&\ast&-2m&\ast&-72o&+218p&-54q&-52r\\
 0=\phantom{]}&-0{,}042&\ast&\ast&\ast&\ast&\ast&\ast&\ast&\ast&\ast&\ast&\ast&\ast&-20n&-52o&-54p&+136q&-10r\\
 0=]&-0{,}021&-55a&\ast&\ast&\ast&\ast&\ast&\ast&\ast&-20i&\ast&-21l&-27m&\ast&-21o&-52p&-10q&+206r.
\end{array}\)%
}
\end{center}
[Die Lösung dieser Normalgleichungen beeinflusst das Tausendstel einer Bogensekunde in den angenommenen Azimuten nicht mehr.\\
Mittlerer Fehler eines einmal gemessenen Winkels \(=\)] \(\pm2''{,}5106.\)
\end{minipage}%
}\hfil\par

% Printed 271; principal PDF 278; Google PDF 282.
\par\clearpage
\noindent\null\hfill{\small\textbf{STATIONSAUSGLEICHUNGEN.}}\hfill\makebox[0pt][r]{\fontsize{13}{16}\selectfont 271}\par\medskip

\vspace*{2.8\baselineskip}
{\centering
[4.]\par\medskip
\{Ausgleichung der auf dem \textls[120]{Windberg} gemessenen Winkel\\
{[}von Art.-Lt. F. \textsc{Hartmann}].\par}
\[
\begin{array}{@{}r@{\enspace}l@{\;}r@{\;}r@{\;}r@{}l@{}}
 &\text{Zielpunkte}&\multicolumn{4}{c@{}}{\text{[Annahme]}}\\
 1&\text{\hbox to 12em{Kirchhesepe\hfil.\hfil.\hfil.\hfil.\hfil.\hfil.\hfil.\hfil.\hfil.\hfil.\hfil.}}&37^\circ&1'&38''{,}0&\\
 2&\text{\hbox to 12em{Kloster ter Apel\hfil.\hfil.\hfil.\hfil.\hfil.\hfil.\hfil.}}&91&23&46{,}2&+a\\
 3&\text{\hbox to 12em{Onstwedde\hfil.\hfil.\hfil.\hfil.\hfil.\hfil.\hfil.\hfil.\hfil.\hfil.\hfil.\hfil.}}&119&22&52{,}6&+b\\
 4&\text{\hbox to 12em{Leer\hfil.\hfil.\hfil.\hfil.\hfil.\hfil.\hfil.\hfil.\hfil.\hfil.\hfil.\hfil.\hfil.\hfil.\hfil.\hfil.\hfil.}}&174&5&30{,}3&+c\\
 5&\text{\hbox to 12em{Bassel\hfil.\hfil.\hfil.\hfil.\hfil.\hfil.\hfil.\hfil.\hfil.\hfil.\hfil.\hfil.\hfil.\hfil.\hfil.\hfil.}}&205&42&28{,}0&+d\\
 6&\text{\hbox to 12em{Molbergen\hfil.\hfil.\hfil.\hfil.\hfil.\hfil.\hfil.\hfil.\hfil.\hfil.\hfil.\hfil.}}&276&35&42{,}5&+e\\
 7&\text{\hbox to 12em{Queckenberg, Signal\hfil.\hfil.\hfil.\hfil.}}&339&44&55{,}9&+f.
\end{array}
\]
\[
\begin{array}{@{}c@{\quad}r@{\;}r@{\;}r@{\;}c@{\;}c@{\;}c@{\;}c@{\;}c@{\;}c@{\;}c@{\quad}r@{\;}r@{\;}r@{}}
 [\text{Winkel}&\multicolumn{3}{c}{\text{Annahme}}&&&&&&&&
 \multicolumn{3}{c@{}}{\text{Beobachtung}] }\\
 1.2&54^\circ&22'&8''{,}2&+a&&&&&&=&54^\circ&22'&7''{,}575\\
 1.3&82&21&14{,}6&\ast&+b&&&&&=&82&21&13{,}925\\
 1.4&137&3&52{,}3&\ast&\ast&+c&&&&=&137&3&52{,}650\\
 1.5&168&40&50{,}0&\ast&\ast&\ast&+d&&&=&168&40&49{,}575\\
 6.1&120&25&55{,}5&\ast&\ast&\ast&\ast&-e&&=&120&25&55{,}400\\
 7.1&57&16&42{,}1&\ast&\ast&\ast&\ast&\ast&-f&=&57&16&41{,}900\\[5pt]
 2.3&27&59&6{,}4&-a&+b&&&&&=&27&59&6{,}237\\
 2.4&82&41&44{,}1&-a&\ast&+c&&&&=&82&41&45{,}100\\
 2.5&114&18&41{,}8&-a&\ast&\ast&+d&&&=&114&18&41{,}465\\
 2.6&185&11&56{,}3&-a&\ast&\ast&\ast&+e&&=&185&11&55{,}500\\
 7.2&111&38&50{,}3&+a&\ast&\ast&\ast&\ast&-f&=&111&38&51{,}240\\[5pt]
 3.4&54&42&37{,}7&\ast&-b&+c&&&&=&54&42&37{,}625\\
 3.5&86&19&35{,}4&\ast&-b&\ast&+d&&&=&86&19&36{,}900\\
 3.6&157&12&49{,}9&\ast&-b&\ast&\ast&+e&&=&157&12&50{,}350\\
 7.3&139&37&56{,}7&\ast&+b&\ast&\ast&\ast&-f&=&139&37&58{,}450
\end{array}
\]
